\section{Viewport dan Unit Responsif}

Viewport adalah area tampilan browser yang terlihat oleh pengguna. Meta tag \code{<meta name="viewport" content="width=device-width, initial-scale=1.0">} wajib ada di setiap halaman responsif---tanpanya, mobile browser akan merender halaman pada lebar desktop (biasanya 980px) lalu memperkecilnya, membuat teks sangat kecil. Unit viewport CSS memungkinkan ukuran elemen relatif terhadap dimensi viewport \cite{mdnCss}.

\begin{table}[htbp]
\centering
\begin{tabular}{|P{1.5cm}|P{4cm}|P{6.5cm}|}
\hline
\textbf{Unit} & \textbf{Relatif Terhadap} & \textbf{Contoh Penggunaan} \\
\hline
\code{vw} & 1\% lebar viewport & Hero section \code{font-size: calc(1rem + 2vw)} \\
\hline
\code{vh} & 1\% tinggi viewport & Full-height hero: \code{height: 100vh} \\
\hline
\code{vmin} & 1\% dimensi terkecil & Layout square responsif \\
\hline
\code{vmax} & 1\% dimensi terbesar & Jarang digunakan \\
\hline
\code{dvh} & Dynamic viewport height & Mengatasi masalah address bar mobile \\
\hline
\code{svh} & Small viewport height & Tinggi tanpa address bar \\
\hline
\end{tabular}
\caption{Unit viewport CSS dan penggunaannya}
\end{table}

Perhatian khusus pada \code{100vh} di perangkat mobile: address bar browser yang muncul/hilang menyebabkan \code{100vh} kadang lebih tinggi dari area yang benar-benar terlihat. Unit modern \code{dvh} (dynamic viewport height) dan \code{svh} (small viewport height) mengatasi masalah ini. Fungsi \code{clamp(min, preferred, max)} sangat berguna untuk \textbf{fluid typography}: \code{font-size: clamp(1rem, 2.5vw, 2rem)} membuat teks yang otomatis menyesuaikan tanpa media query \cite{webdevCss}.

\begin{konsep}
\textbf{Fluid Typography dengan clamp().} Alih-alih mendefinisikan \code{font-size} berbeda di setiap breakpoint, gunakan \code{clamp()} untuk tipografi yang mulus: \code{font-size: clamp(1rem, 1rem + 1vw, 2.5rem);} --- teks tumbuh secara proporsional dari minimum 1rem ke maksimum 2.5rem berdasarkan lebar viewport. Ini mengurangi jumlah media queries dan menghasilkan transisi ukuran teks yang lebih halus \cite{cssTricks}.
\end{konsep}

